\documentclass[10pt,a4paper]{amsart}
\usepackage[margin=19mm]{geometry}
\usepackage{amsmath,amssymb,mathtools}
\usepackage[scr=rsfs,cal=euler]{mathalfa}
\usepackage{xfrac}
\usepackage[unicode,hidelinks,bookmarks=false]{hyperref}
\newcommand{\ie}{i.e.\ }
\newcommand{\cf}{cf.\ }
\newcommand{\resp}{respectively\ }
\newcommand{\Subsection}{\subsection}
\providecommand{\nobreakdash}{\nobreak\hbox{-}}
\setlength{\emergencystretch}{2em}
\multlinegap0pt
\usepackage{amssymb}
\newtheorem{lemma}{Lemma}
\newtheorem{proposition}{Proposition}
\newcommand{\R}{\mathbb{R}}
\title{Self-similar solutions of the three-dimensional Muskat problem with surface tension}
\author{Lizhe Wan, Jiaqi Yang}
\date{Source: arXiv:2608.00430v1 (CC BY 4.0)}
\begin{document}
\maketitle

\noindent\emph{Source and licence.} Self-similar solutions of the three-dimensional Muskat problem with surface tension. Version arXiv:2608.00430v1 (CC BY 4.0), distributed by arXiv under the Creative Commons Attribution 4.0 International licence, http://creativecommons.org/licenses/by/4.0/, as recorded in the arXiv licence metadata for this exact version and shown on the abstract page https://arxiv.org/abs/2608.00430v1. Full licence notice: \texttt{licenses/arxiv-2608.00430-CC-BY-4.0.txt}.\par\smallskip

\noindent\textit{Editorial scope.} Counted unit: the article's proposition prop:finite-order-DN (source0.tex lines 1613--1632). The article's complete original proof of it is delivered with it (source0.tex lines 1634--1855, one \texttt{proof} environment of 222 lines). No mathematics is written, repaired or re-derived: the statement and the proof are the article's own lines, in the article's own notation, and no retained line is substituted or re-flowed.  Retained uncounted as the source-local context the proof needs: lines 1289--1314; lines 1355--1397; lines 1442--1454; lines 1535--1556.  Nothing was omitted from any retained span, so the package carries no omission record.  No retained line contains a citation, so the wrapper carries no bibliography and no bibliography entry.  No retained line contains a figure environment, an image-inclusion command or drawing code, so no image is delivered and the package declares no asset dependency.  Editorial formatting only, in addition: an AMS \texttt{amsart} preamble that restates the article's own declarations and macros used by the retained lines, namely the environments \texttt{lemma}, \texttt{proposition} on separate counters, together with the standard packages those lines need.  The retained environments print in this document as Lemma 1, Proposition 1 in order of appearance, whereas the article prints the same items with the numbering of its own full document.  The complete native source of the article as distributed in the arXiv e-print for this exact version is bundled unchanged as \texttt{sources/206540/source0.tex}.
\begin{lemma}[Coefficient algebra and composition]
\label{lem:coefficient-algebra}
Let \(r>1\).
Then \(\mathcal E_r\) is a Banach algebra and
\begin{equation*}
 \|AB\|_{\mathcal E_r} \leq C_r \left( \|A\|_{L^\infty}\|B\|_{\mathcal E_r} + \|B\|_{L^\infty}\|A\|_{\mathcal E_r} \right).
\end{equation*}

Let \(\Phi\) be a \(C^\infty\) function on a neighborhood of the closed ball \(\{p:|p|\leq M\}\) in a finite-dimensional vector space, and define the Nemytskii map
\[
 \mathcal N_\Phi(A):=\Phi\circ A.
\]
Then \(\mathcal N_\Phi\) is \(C^4\) from 
\[
 \left\{ A\in\mathcal E_r: \|A\|_{L^\infty}<M \right\}
\]
to \(\mathcal E_r\).
More precisely, on every set on which
\[
 \|A\|_{\mathcal E_r}\leq M_1, \qquad \|A\|_{L^\infty}\leq M_0<M,
\]
and for \(1\leq m\leq4\), one has
\begin{equation*}
 \left\| D^m\mathcal N_\Phi(A)[B_1,\ldots,B_m] \right\|_{\mathcal E_r} \leq C_{\Phi,M_0,M_1,r,m} \prod_{\nu=1}^m \|B_\nu\|_{\mathcal E_r}.
\end{equation*}
\end{lemma}

\begin{lemma}[Regularized extension bounds]
\label{lem:regularized-extension-two-level}
Fix an even function \(\chi\) such that
\[
 \chi(x)\in C^\infty_c(\R^2), \qquad \int_{\R^2}\chi(x) dx =1,
\]
and define its rescaling by
\[
 \chi_\tau(x):=\tau^{-2}\chi(x/\tau), \qquad \tau>0.
\]
For \(h\in\mathfrak A_r\) and \(z<0\), we define
\[
 H_h(x,z):=\chi_{-z}*h(x), \qquad E_h:=\nabla_{x,z}H_h.
\]
We also set \(H_h(x,0)=h(x)\) in the trace sense.
Then
\begin{equation}
\label{eq:Eh-Linfty}
 \|E_h\|_{L^\infty_{x,z}} \leq C_\chi\|\nabla h\|_{L^\infty}.
\end{equation}
Moreover, for every \(N\geq0\),
\begin{equation*}
 \|\dot\Delta_jE_h(\cdot,z)\|_{L^2} \leq C_{N,\chi} 2^j(1+2^j|z|)^{-N} \|\dot\Delta_jh\|_{L^2}.
\end{equation*}
Consequently,
\begin{align}
\label{eq:Eh-high}
 \left( \sum_{j\in\mathbb Z} 2^{2rj} \|\dot\Delta_jE_h\|_{L^\infty_zL^2_x}^2 \right)^{1/2}+ \left( \sum_{j\in\mathbb Z} 2^{(2r+1)j} \|\dot\Delta_jE_h\|_{L^2_{x,z}}^2 \right)^{1/2} \leq C_{r,\chi}\|h\|_{\dot H^{r+1}}.
\end{align}
In particular,
\begin{equation*}
 \|E_h\|_{\mathcal E_r} \leq C_{r,\chi}\mathcal A_r(h).
\end{equation*}
Finally, as \(z\uparrow0\),
\begin{equation}
\label{eq:regularized-extension-boundary-traces}
\begin{aligned}
 H_h(\cdot,z)&\longrightarrow h \qquad \text{locally uniformly and in }\dot H^{r+1},\\
 \nabla_xH_h(\cdot,z)&\longrightarrow\nabla h \quad\text{in }\dot H^r,
 \qquad \partial_zH_h(\cdot,z)\longrightarrow0 \quad\text{in }\dot H^r.
\end{aligned}
\end{equation}
\end{lemma}

\begin{lemma}[Flat mixed first-order estimate]
\label{lem:flat-mixed-elliptic}
Let \(f\) satisfy \(\mathcal D_r(f)<\infty\), and let \(F_0,F_1\in\mathcal Y_r\).
Consider
\begin{align*}
 (\partial_z-|D|)V&=W+F_0, &V|_{z=0}&=f,\\
 (\partial_z+|D|)W&=|D|F_1, &\lim_{z\to-\infty}W(z)&=0.
\end{align*}
Then
\begin{equation*}
 \|V\|_{\mathcal X_r} + \|W\|_{\mathcal Y_r} + \|W|_{z=0}\|_{\dot H^r} \leq C_r \left( \mathcal D_r(f) + \|F_0\|_{\mathcal Y_r} + \|F_1\|_{\mathcal Y_r} \right).
\end{equation*}
\end{lemma}

\begin{lemma}[Mixed tame product estimate]
\label{lem:mixed-tame-product}
Let \(r>1\).
Suppose that \(a\in L^\infty(\R^3_-)\) and
\[
 \|a\|_{\mathcal E_r^+} := \left( \sum_{j>0} 2^{2rj} \|\dot\Delta_ja\|_{L^\infty_zL^2_x}^2 \right)^{1/2} + \left( \sum_{j>0} 2^{(2r+1)j} \|\dot\Delta_ja\|_{L^2_{x,z}}^2 \right)^{1/2} <\infty.
\]
If \(B=\partial_\ell V\), where
\(\partial_\ell\in
\{\partial_{x_1},\partial_{x_2},\partial_z\}\), then
\begin{equation}
\label{eq:mixed-tame-product}
 \|aB\|_{\mathcal Y_r} \leq C_r \left( \|a\|_{L^\infty}\|V\|_{\mathcal X_r} + \|a\|_{\mathcal E_r^+} \|\nabla_{x,z}V\|_{L^\infty_{x,z}} \right).
\end{equation}
Moreover,
\begin{equation}
\label{eq:X-to-Linfty}
 \|\nabla_{x,z}V\|_{L^\infty_{x,z}} \leq C_r\|V\|_{\mathcal X_r}.
\end{equation}

If \(a\) is a finite product or a smooth composition of coefficients belonging to \(\mathcal E_r\), then \eqref{eq:mixed-tame-product} remains valid after estimating \(\|a\|_{L^\infty}+\|a\|_{\mathcal E_r^+}\) by Lemma~\ref{lem:coefficient-algebra}.
\end{lemma}

\begin{proposition}[Finite-order cone-adapted DN estimate]
\label{prop:finite-order-DN}
Let \(3/2<r<2\).
There exists \(\delta_r>0\) such that the map
\[
 \eta\longmapsto G(\eta)
\]
is \(C^3\) from the open ball
\[
 \{\eta\in\mathfrak A_r:\mathcal A_r(\eta)<2\delta_r\}
\]
into \(\mathcal L(\mathfrak D_r,\dot H^r(\R^2))\).
More precisely, whenever \(\mathcal A_r(\eta)\leq\delta_r\) and
\(k=0,1,2,3\),
\begin{equation}
\label{eq:finite-order-DN}
 \left\| D_\eta^kG(\eta)[h_1,\ldots,h_k]f \right\|_{\dot H^r} \leq C_{r,k}\, \mathcal D_r(f) \prod_{\nu=1}^k\mathcal A_r(h_\nu).
\end{equation}
For \(k=0\), the product on the right-hand side is understood to be one, and the constants are uniform under the stated smallness condition.
\end{proposition}

\begin{proof}
We first work with smooth finite-frequency data; the estimates are uniform
and therefore extend to the completed spaces. Set
\[
 \rho_\eta(x,z)=z+H_\eta(x,z),
 \qquad v_\eta(x,z)=\phi_f(x,\rho_\eta(x,z)).
\]
For \(\delta_r\) small enough, Lemma~\ref{lem:regularized-extension-two-level} gives
\[
\tfrac12\leq\partial_z\rho_\eta\leq\tfrac32, \qquad |H_\eta(x,z)-\eta(x)|
 \lesssim |z|\|\nabla\eta\|_{L^\infty}.
\]
Together with the boundary traces, these estimates show that \((x,z)\mapsto(x,\rho_\eta(x,z))\) is a bi-Lipschitz map from the lower
half-space onto \(\Omega_\eta\).

If we write \(E_\eta:=\nabla_{x,z}H_\eta\), a direct change-of-variables gives
\begin{equation*}
 \operatorname{div}_{x,z} \left( \mathsf A_\eta\nabla_{x,z}v_\eta \right)=0, \qquad
 \mathsf A_\eta  :=
 \begin{pmatrix}
 (1+\partial_zH_\eta)I_2 & -\nabla H_\eta\\[2mm]
 -(\nabla H_\eta)^T & \displaystyle \frac{1+|\nabla H_\eta|^2} {1+\partial_zH_\eta}
 \end{pmatrix}.
\end{equation*}
The matrix \(\mathsf A_\eta\) is uniformly elliptic.
Indeed, if \(a=\nabla H_\eta\), \(b=\partial_zH_\eta\), and \((\xi,\zeta)\in\R^2\times\R\), then
\[
 \mathsf A_\eta \binom{\xi}{\zeta} \cdot \binom{\xi}{\zeta} = \frac1{1+b} \left( |(1+b)\xi-a\zeta|^2+\zeta^2 \right),
\]
and \(a,b\) are uniformly small.

Expanding the divergence-form equation gives the equivalent system
\begin{equation*}
 \Delta_{x,z}v_\eta = \partial_zQ_a(\eta,v_\eta) + |D|Q_b(\eta,v_\eta), \qquad v_\eta|_{z=0}=f,
\end{equation*}
where
\begin{align*}
 Q_a(\eta,V) &:= \nabla H_\eta\cdot\nabla_xV
 -\frac{|\nabla H_\eta|^2-\partial_zH_\eta}
 {1+\partial_zH_\eta}\,\partial_zV,\\
 Q_b(\eta,V) &:= |D|^{-1}\operatorname{div}_x
 \left(  \nabla H_\eta\,\partial_zV
  -\partial_zH_\eta\,\nabla_xV \right).
\end{align*}

Define
\[
 w_\eta := (\partial_z-|D|)v_\eta - Q_a(\eta,v_\eta).
\]
Using the decomposition of the Laplacian
\[
 \Delta_{x,z} = (\partial_z+|D|)(\partial_z-|D|),
\]
we obtain
\begin{equation*}
 (\partial_z+|D|)w_\eta = |D| \left( Q_b(\eta,v_\eta)-Q_a(\eta,v_\eta) \right).
\end{equation*}

At \(z=0\), put \(a=\nabla_xH_\eta=\nabla\eta\) and
\(b=\partial_zH_\eta\).  The chain rule gives
\[
 G(\eta)f
 =-a\cdot\nabla_xv_\eta+\frac{1+|a|^2}{1+b}\partial_zv_\eta
 =\bigl(\partial_zv_\eta-Q_a(\eta,v_\eta)\bigr)\big|_{z=0}.
\]
Since \(v_\eta(0)=f\),
\begin{equation}
\label{eq:DN-trace-first-order}
 G(\eta)f=|D|f+w_\eta(0).
\end{equation}

Writing \(R_j:=|D|^{-1}\partial_{x_j}\), we have
\[
 Q_b(\eta,V) = \sum_{j=1}^2 R_j
 \left(  \partial_{x_j}H_\eta\,\partial_zV
  -\partial_zH_\eta\,\partial_{x_j}V
 \right).
\]
Thus \(Q_a\), and the expressions inside the Riesz transforms in
\(Q_b\), are sums of products \(C(E_\eta)\partial_\ell V\), where
\(C\) is smooth and \(C(0)=0\). 
Lemma~\ref{lem:coefficient-algebra} and
Lemma~\ref{lem:regularized-extension-two-level} give
\[
 \|C(E_\eta)\|_{L^\infty} + \|C(E_\eta)\|_{\mathcal E_r^+} \lesssim_r \mathcal A_r(\eta).
\]
Lemma~\ref{lem:mixed-tame-product} and the \(L^2\)-boundedness of the horizontal Riesz transforms therefore yield
\begin{equation}\label{eq:Q-base-bound}
 \|Q_a(\eta,V)\|_{\mathcal Y_r}
 + \|Q_b(\eta,V)\|_{\mathcal Y_r}
 \lesssim_r \mathcal A_r(\eta)\|V\|_{\mathcal X_r}.
\end{equation}

For \(V\in\mathcal X_r\), define
\begin{equation*}
 \mathcal W_\eta[V](z) := \int_{-\infty}^z e^{-(z-z')|D|}|D| \left( Q_b(\eta,V)-Q_a(\eta,V) \right)(z')\,dz'
\end{equation*}
and
\begin{equation*}
 \mathcal K_\eta V(z) := \int_0^z e^{(z-z')|D|} \left( \mathcal W_\eta[V](z') + Q_a(\eta,V)(z') \right)\,dz'.
\end{equation*}
Applying Lemma~\ref{lem:flat-mixed-elliptic} with zero boundary datum,
\(F_0=Q_a(\eta,V)\), and \(F_1=Q_b(\eta,V)-Q_a(\eta,V)\), and then using
\eqref{eq:Q-base-bound}, we obtain
\begin{equation} \label{eq:K-W-base-bound}
\begin{aligned}
 \|\mathcal K_\eta V\|_{\mathcal X_r} + \|\mathcal W_\eta[V]\|_{\mathcal Y_r} + \|\mathcal W_\eta[V](0)\|_{\dot H^r}\leq C_r\mathcal A_r(\eta) \|V\|_{\mathcal X_r}.
\end{aligned}
\end{equation}

The flattened extension satisfies
\begin{equation*}
 v_\eta = P[f]+\mathcal K_\eta v_\eta, \qquad P[f](z):=e^{z|D|}f.
\end{equation*}
Choose \(\delta_r\) so that \(\|\mathcal K_\eta\| \leq\frac12\), then $(I-\mathcal K_\eta)$ is invertible and
\[
 v_\eta = (I-\mathcal K_\eta)^{-1}P[f].
\]
The Poisson estimate and \eqref{eq:K-W-base-bound} give
\begin{equation*}
 \|v_\eta\|_{\mathcal X_r} + \|w_\eta\|_{\mathcal Y_r} + \|w_\eta(0)\|_{\dot H^r} \leq C_r\mathcal D_r(f).
\end{equation*}
The trace formula \eqref{eq:DN-trace-first-order} proves
\eqref{eq:finite-order-DN} for \(k=0\).

The map \(h\mapsto E_h\) is linear from \(\mathfrak A_r\) to \(\mathcal E_r\).  
Since \(1+\partial_zH_\eta\geq\frac12\), the coefficient algebra, composition, and mixed product estimates give, for
\(1\leq m\leq4\),
\begin{equation}
\label{eq:Q-shape-operator-bound}
 \sum_{i = a, b}\left\| D_\eta^mQ_i(\eta,V) [h_1,\ldots,h_m] \right\|_{\mathcal Y_r}\leq C_{r,m} \left( \prod_{\nu=1}^m \mathcal A_r(h_\nu) \right) \|V\|_{\mathcal X_r}.
\end{equation}
Each differentiated term consists of \(E_{h_1},\ldots,E_{h_m}\), a smooth bounded function of \(E_\eta\), and one component of \(\nabla V\);
the tame estimates place at most one factor in a homogeneous norm.

Define \(\mathcal Z_\eta V:=\mathcal W_\eta[V](0)\).
Because the two Duhamel operators in Lemma~\ref{lem:flat-mixed-elliptic} are bounded and independent of \(\eta\), differentiating them under the integral sign and using
\eqref{eq:Q-shape-operator-bound} shows that \(\mathcal K_\eta\) and \(\mathcal Z_\eta\) are \(C^3\) in operator norm.  
More precisely, for
\(1\leq m\leq3\),
\begin{equation}
\label{eq:K-shape-bound}
\begin{aligned}
 \left\| D_\eta^m\mathcal K_\eta [h_1,\ldots,h_m]V \right\|_{\mathcal X_r}+ \left\| D_\eta^m\mathcal Z_\eta [h_1,\ldots,h_m]V \right\|_{\dot H^r}\leq C_{r,m} \left( \prod_{\nu=1}^m \mathcal A_r(h_\nu) \right) \|V\|_{\mathcal X_r}.
\end{aligned}
\end{equation}
For example, \(D_\eta^m\mathcal W_\eta[V]\) is obtained by replacing \(Q_b-Q_a\) in its defining integral by \(D_\eta^m(Q_b-Q_a)[h_1,\ldots,h_m]V\), and the formula for
\(D_\eta^m\mathcal K_\eta\) has the analogous two terms.  
The estimate with \(m=4\) following from \eqref{eq:Q-shape-operator-bound} controls the third-order Taylor remainder and proves continuity of the third derivatives.

Set
\[
 \mathcal R_\eta:=(I-\mathcal K_\eta)^{-1}.
\]
By \eqref{eq:K-W-base-bound} and the choice of \(\delta_r\),
\(\|\mathcal K_\eta\|_{\mathcal L(\mathcal X_r)} \leq\frac12\).
Hence the Neumann series gives
\[
 \|\mathcal R_\eta\|_{\mathcal L(\mathcal X_r)}\leq2.
\]
Since \(\eta\mapsto\mathcal K_\eta\) is \(C^3\) in operator norm, so
is \(\eta\mapsto\mathcal R_\eta\). Differentiating
\((I-\mathcal K_\eta)\mathcal R_\eta=I\) gives
\[
 D_\eta\mathcal R_\eta[h]
 = \mathcal R_\eta
 \bigl(D_\eta\mathcal K_\eta[h]\bigr) \mathcal R_\eta.
\]
Repeated differentiation shows that \(D_\eta^k\mathcal R_\eta\), \(k\leq3\), is a finite sum of terms alternating factors \(\mathcal R_\eta\) with derivatives
\(D_\eta^m\mathcal K_\eta\), where the derivative directions are partitioned among the latter factors. 
Hence \eqref{eq:K-shape-bound}
gives
\[
 \left\|D_\eta^k\mathcal R_\eta[h_1,\ldots,h_k]\right\|_{\mathcal L(\mathcal X_r)}
 \leq C_{r,k}\prod_{\nu=1}^k\mathcal A_r(h_\nu).
\]
Since \(v_\eta=\mathcal R_\eta P[f]\), the Poisson estimate yields, for \(1\leq k\leq3\),
\begin{equation*}
 \left\| D_\eta^kv_\eta [h_1,\ldots,h_k] \right\|_{\mathcal X_r} \leq C_{r,k} \mathcal D_r(f) \prod_{\nu=1}^k \mathcal A_r(h_\nu).
\end{equation*}

The trace identity reads \(G(\eta)f=|D|f+\mathcal Z_\eta\mathcal R_\eta P[f]\).  Leibniz rule, \eqref{eq:K-shape-bound}, and the preceding resolvent bounds give
\[
 \left\| D_\eta^kG(\eta) [h_1,\ldots,h_k]f \right\|_{\dot H^r} \leq C_{r,k} \mathcal D_r(f) \prod_{\nu=1}^k \mathcal A_r(h_\nu), \qquad 1\leq k\leq3.
\]
The same operator formulas are meaningful directly for every
\(\eta\in\mathfrak A_r\) in the small ball and every
\(f\in\mathfrak D_r\); no density in the Lipschitz norm is required.
The coefficient estimates define \(Q_a,Q_b\), hence
\(\mathcal K_\eta\), on \(\mathcal X_r\), and the contraction formula
defines \(v_\eta\).  Substitution into the two first-order equations shows
in distributions that
\[
 \operatorname{div}_{x,z}(\mathsf A_\eta\nabla_{x,z}v_\eta)=0,
 \qquad v_\eta|_{z=0}=f.
\]
Moreover, \(\nabla v_\eta\in L^2(\R^3_-)\), so this is an energy
solution.  Composing with the inverse bi-Lipschitz flattening map gives an
energy harmonic extension in \(\Omega_\eta\), with boundary trace \(f\).
It is unique: the difference
\(u\in\dot H^1_0(\R^3_-)\) of two energy solutions satisfies
\[
 \int_{\R^3_-}\mathsf A_\eta\nabla u\cdot\nabla u=0.
\]
Ellipticity and the zero trace imply \(u=0\).  
Moreover,
\[
 (\mathsf A_\eta\nabla v_\eta)_3
 =\partial_zv_\eta-Q_a(\eta,v_\eta)=|D|v_\eta+w_\eta
 \longrightarrow |D|f+w_\eta(0)
\]
in boundary distributions.  Here the zero-boundary Duhamel term in
\(v_\eta-P[f]\) tends to zero, while the boundary estimate in
Lemma~\ref{lem:flat-mixed-elliptic} supplies the trace of \(w_\eta\).
Testing the weak equation against an extension \(\Psi\) of \(\psi\in C_c^\infty(\R^2)\) gives
\[
 \int_{\R^3_-}\mathsf A_\eta\nabla v_\eta\cdot\nabla\Psi =\langle |D|f+w_\eta(0),\psi\rangle.
\]
Thus \(|D|f+\mathcal Z_\eta v_\eta\) is the variational
Dirichlet--Neumann trace. This identifies the completed-space construction with the classical operator for smooth data and completes the proof of
\eqref{eq:finite-order-DN}.
\end{proof}

\end{document}
